% Part: computability
% Chapter: recursive-functions
% Section: sequences

\documentclass[../../../include/open-logic-section]{subfiles}

\begin{document}

\olfileid{cmp}{rec}{seq}
\olsection{क्रमिका}

आदिम पुनरावर्ती फलनांचा संच उल्लेखनीयरीत्या सुदृढ आहे. तरीही
\emph{क्रमिका} हाताळण्याची पुरेशी पद्धत विकसित केल्यावर आपल्याला आणखी बरेच
काही करता येईल. नैसर्गिक संख्यांच्या सांत क्रमिका आपण पुढील प्रकारे नैसर्गिक
संख्यांशी ओळखू: $\langle a_0, a_1, a_2, \dots, a_k \rangle$ ही क्रमिका
पुढील संख्येशी संगत असेल:
\[
p_0^{a_0+1} \cdot p_1^{a_1+1} \cdot p_2^{a_2+1} \cdot \dots \cdot
p_k^{a_k+1}.
\]
उदाहरणार्थ, $\langle 2, 7, 3\rangle$ आणि $\langle 2, 7, 3, 0, 0 \rangle$
या क्रमिकांचे संख्यात्मक संकेतांक वेगळे राहतील याची खात्री करण्यासाठी आपण
घातांकांमध्ये एक मिळवतो. रिकाम्या क्रमिकेचा संकेतांक म्हणून $0$ आणि~$1$
दोन्ही घेता येतात; निश्चिततेसाठी $\emptyseq$ हे~$0$ दर्शवू दे.

क्रमिकांचे हे सांकेतीकरण चालते याचे कारण तथाकथित अंकगणिताचे मूलभूत प्रमेय
होय: प्रत्येक $n \ge 2$ नैसर्गिक संख्या पुढील रूपात एकमेव रीतीने लिहिता
येते:
\[
n = p_0^{a_0} \cdot p_1^{a_1} \cdot \dots \cdot p_k^{a_k}
\]
येथे $a_k \ge 1$. त्यामुळे $\tuple{}(a_0,
\dots, a_k) = \tuple{a_0, \dots, a_k}$ हे चित्रण एकास-एक आहे: वेगवेगळ्या
क्रमिका वेगवेगळ्या संख्यांकडे पाठवल्या जातात आणि प्रत्येक संख्येशी जास्तीत
जास्त एकच क्रमिका संगत असते.

आता क्रमिकेची लांबी ठरवणे, तिचा $i$ वा घटक ठरवणे, क्रमिकेच्या शेवटी एक
घटक जोडणे आणि दोन क्रमिकांची जोडणी करणे या सर्व क्रिया आदिम पुनरावर्ती
आहेत हे आपण दाखवू.

\begin{prop}
  क्रमिका~$s$ ची लांबी परत करणारे $\len{s}$ हे फलन आदिम पुनरावर्ती आहे.
\end{prop}

\begin{proof}
  $R(i, s)$ हा संबंध पुढीलप्रमाणे परिभाषित करू:
  \[
  R(i, s) \text{ तेव्हा आणि केवळ तेव्हाच }
  p_i \mid s \land p_{i+1} \nmid s.
  \]
  $R$ हा संबंध स्पष्टपणे आदिम पुनरावर्ती आहे. $s$ हा रिकाम्या नसलेल्या
  क्रमिकेचा संकेतांक असेल, म्हणजे
  \[
  s = p_0^{a_0+1} \cdot \dots \cdot p_{k}^{a_{k}+1},
  \]
  तेव्हा $p_i \mid s$ असणारी $p_i$ ही सर्वांत मोठी अविभाज्य संख्या असल्यास,
  म्हणजे $i = k$ असल्यास, $R(i,s)$ सत्य असते. म्हणून $s$ ची लांबी $i+1$
  असते तेव्हा आणि केवळ तेव्हाच $p_i$ ही~$s$ ला भाग जाणारी सर्वांत मोठी
  अविभाज्य संख्या असते. त्यामुळे पुढील व्याख्या करता येते:
  \[
  \len{s} =
  \begin{cases}
    0 & \text{जर $s = 0$ किंवा $s = 1$ असेल} \\
    1 + \bmin{i < s}{R(i, s)} & \text{अन्यथा.}
  \end{cases}
  \]
  येथे परिबद्ध लघुतमीकरण वापरता येते, कारण $s$ हा क्रमिकेचा संकेतांक असेल
  तेव्हा $R(i,s)$ पूर्ण करणारा फक्त एकच~$i$ असतो आणि असा $i$ अस्तित्वात
  असल्यास तो~$s$ पेक्षा लहान असतो.
\end{proof}

\begin{prop}
  क्रमिका~$s$ च्या शेवटी $a$ जोडल्यावर मिळणारी क्रमिका परत करणारे
  $\fn{append}(s,a)$ हे फलन आदिम पुनरावर्ती आहे.
\end{prop}

\begin{proof}
  $\fn{append}$ ची व्याख्या पुढीलप्रमाणे करता येते:
  \[
  \fn{append}(s,a) =
  \begin{cases}
    2^{a+1} & \text{जर $s = 0$ किंवा $s = 1$ असेल} \\
    s \cdot p_{\len{s}}^{a+1} & \text{अन्यथा.}
  \end{cases}
  \]
\end{proof}

\begin{prop}
  $s$ चा $i$ वा घटक परत करणारे $\fn{element}(s,i)$ हे फलन आदिम
  पुनरावर्ती आहे; येथे पहिल्या घटकाला $0$ वा म्हणतात आणि $i$ हा $s$ च्या
  लांबीपेक्षा मोठा किंवा समान असल्यास फलन $0$ परत करते.
\end{prop}

\begin{proof}
$a$ हा~$s$ चा $i$ वा घटक असतो तेव्हा आणि केवळ तेव्हाच $p_i^{a+1}$ ही
$s$ ला भाग जाणारी~$p_i$ ची सर्वांत मोठी घातशक्ती असते, म्हणजे
$p_i^{a+1} \mid s$ पण $p_i^{a+2} \nmid s$. म्हणून:
  \[
  \fn{element}(s,i) =
  \begin{cases}
    0 & \mbox{जर $i \geq \len{s}$ असेल} \\
    \bmin{a < s}{(p_i^{a+2} \nmid s)} & \text{अन्यथा.}
  \end{cases}
  \]
\end{proof}

वर परिभाषित केलेल्या फलनांची पूर्ण नावे वापरण्याऐवजी आपण अधिक संक्षिप्त
चिन्हांकन आणू. $\fn{element}(s,i)$ ऐवजी $(s)_i$ वापरू आणि
$\tuple{s_0, \dots, s_k}$ हे पुढील अभिव्यक्तीचे संक्षिप्त रूप मानू:
\[
\fn{append}(\fn{append}(\dots \fn{append}(\emptyseq,s_0)
\dots),s_k).
\]
$s$ ची लांबी~$k$ असेल, तर $s$ चे घटक $(s)_0$, \dots,~$(s)_{k-1}$
असतात, हे लक्षात घ्या.

\begin{prop}
दोन क्रमिकांची जोडणी करणारे $\fn{concat}(s,t)$ हे फलन आदिम पुनरावर्ती
आहे.
\end{prop}

\begin{proof}
  आपल्याला पुढील गुणधर्म असलेले $\fn{concat}$ हे फलन हवे आहे:
  \[
    \fn{concat}(\tuple{a_0, \dots, a_k}, \tuple{b_0, \dots, b_l}) =
    \tuple{a_0, \dots, a_k, b_0, \dots, b_l}.
  \]
  $t$ चे पहिले $n$ घटक~$s$ च्या शेवटी जोडणारे
  $\fn{hconcat}(s,t,n)$ हे ``साहाय्यक'' फलन वापरू. त्याची आदिम
  पुनरावर्तनाने पुढीलप्रमाणे व्याख्या करता येते:
  \begin{align*}
    \fn{hconcat}(s,t,0) & = s\\
    \fn{hconcat}(s,t,n+1) & = \fn{append}(\fn{hconcat}(s,t,n),(t)_n)
    \intertext{मग $\fn{concat}$ ची व्याख्या पुढीलप्रमाणे करता येते:}
    \fn{concat}(s,t) & = \fn{hconcat}(s,t,\len{t}).
  \end{align*}
\end{proof}

$\fn{concat}(s,t)$ ऐवजी आपण $s \concat t$ असे लिहू.

क्रमिकेच्या लांबीच्या आणि तिच्या सर्वांत मोठ्या घटकाच्या साहाय्याने तिच्या
संख्यात्मक संकेतांकावर परिबंध घालता येणे आपल्याला उपयोगी पडेल. $s$ ही
लांबी~$k$ असलेली क्रमिका आहे आणि तिचा प्रत्येक घटक एखाद्या~$x$ पेक्षा
लहान किंवा समान आहे असे समजा. मग $s$ चे जास्तीत जास्त $k$ अविभाज्य अवयव
असतात; त्यांपैकी प्रत्येक जास्तीत जास्त~$p_{k-1}$ असतो आणि~$s$ च्या
अविभाज्य अवयवीकरणात प्रत्येकाचा घातांक जास्तीत जास्त $x+1$ असतो. दुसऱ्या
शब्दांत, धन लांबीसाठी आपण
\[
\fn{sequenceBound}(x,k) = p_{k-1}^{k \cdot (x+1)},
\]
अशी व्याख्या केली, तर वर वर्णन केलेल्या क्रमिका~$s$ चा संख्यात्मक संकेतांक
जास्तीत जास्त~$\fn{sequenceBound}(x,k)$ असतो. शून्य लांबीच्या क्रमिकेसाठी
क्रमिका-परिबंधाचे मूल्य एक मानू.

क्रमिकांवरील असा परिबंध आपल्याला परिबद्ध शोध वापरून नवीन फलने परिभाषित
करण्याची पद्धत देतो. उदाहरणार्थ, परिबद्ध शोध वापरून $\fn{concat}$ ची
व्याख्या करता येते. आपण शोधत असलेल्या वस्तूचे, म्हणजे जोडलेल्या क्रमिकेच्या
संख्यात्मक संकेतांकाचे, आदिम पुनरावर्ती \emph{विनिर्देश} आणि शोध किती दूरवर
करायचा त्यावरील परिबंध एवढेच लिहावे लागते. पुढील व्याख्या हे काम करते:
\begin{align*}
  \fn{concat}(s,t) = {} & \bmin{v < \fn{sequenceBound}(s+t,\len{s} +
    \len{t})}{} \\
  & \quad(\len{v} = \len{s} + \len{t} \land {}\\
    & \qquad \bforall{i < \len{s}}{((v)_i = (s)_i) \land {} \\
      & \qquad \bforall{j < \len{t}}{((v)_{\len{s}+j} = (t)_j)})}
\end{align*}

\begin{prob}
पुढील गुणधर्म असलेले आदिम पुनरावर्ती फलन~$\fn{sconcat}(s)$ अस्तित्वात आहे
हे दाखवा:
\[
\fn{sconcat}(\tuple{s_0, \dots, s_k}) = s_0 \concat \dots \concat s_k.
\]
\end{prob}

\begin{prob}
पुढील गुणधर्म असलेले आदिम पुनरावर्ती फलन~$\fn{tail}(s)$ अस्तित्वात आहे
हे दाखवा:
\begin{align*}
  \fn{tail}(\emptyseq) & = 0 \text{ आणि}\\
  \fn{tail}(\tuple{s_0, \dots, s_{k}}) & = \tuple{s_1, \dots, s_{k}}.
\end{align*}
\end{prob}

\begin{prop}
  \ollabel{prop:subseq}
  $s$ च्या $i$ व्या घटकापासून सुरू होणारी आणि लांबी~$n$ असलेली उपक्रमिका
  परत करणारे $\fn{subseq}(s, i, n)$ हे फलन आदिम पुनरावर्ती आहे.
\end{prop}

\begin{proof}
  सराव.
\end{proof}

\begin{prob}
  \olref[cmp][rec][seq]{prop:subseq} सिद्ध करा.
\end{prob}
  
\end{document}
